% ============================================================================
% APÊNDICE B — Soluções dos exercícios
% ============================================================================
\chapter{Soluções dos Exercícios}
\label{apendice:solucoes}

\begin{caixaconceito}
Política editorial: as soluções completas, executáveis e testadas vivem no
repositório de código (\texttt{codigo/--solucoes/}), organizadas por capítulo.
Este apêndice explica o padrão e traz as soluções comentadas dos exercícios
essenciais das Partes I e II.
\end{caixaconceito}

\section{Como usar as soluções}

\begin{itemize}[leftmargin=*, itemsep=0.4em]
  \item \textbf{Primeiro tente sozinho}: a solução só ensina depois que você
  tentou (e errou) — o erro é parte do aprendizado;
  \item \textbf{Compare, não copie}: leia a solução, feche o arquivo e
  reescreva do seu jeito;
  \item \textbf{Teste sempre}: toda solução tem pelo menos um teste
  (Vitest) — se não passa, não é solução;
  \item \textbf{Critérios de avaliação}: os projetos têm checklist no README
  de cada capítulo (você pode usar como rubrica de autoavaliação).
\end{itemize}

\section{Estrutura do diretório de soluções}

\begin{lstlisting}[style=livro, language=Bash, caption={Organização das soluções.}, label={list:sol-estrutura}]
codigo/--solucoes/
  README.md                 # como usar + cobertura
  cap01/  solucoes-cap01.md # gabarito comentado
  cap02/  solucoes-cap02.md
  ...
  cap25/  solucoes-cap25.md
\end{lstlisting}

\section{Soluções comentadas — Capítulo 1}

\begin{exercicios}
  \item \textbf{Frontend vs backend vs full stack}: o frontend roda no
  navegador (HTML/CSS/JS); o backend roda no servidor (APIs, regras de
  negócio); o full stack transita pelas duas camadas, além de banco e infra.
  \item \textbf{Top 3 tecnologias} (SO Survey 2025): JavaScript (linguagem),
  SQL (tecnologia geral), React (framework web) — com PostgreSQL como banco nº 1
  e Node.js líder entre runtimes.
  \item \textbf{Jornada da requisição}: DNS $\rightarrow$ conexão TLS
  $\rightarrow$ requisição HTTP $\rightarrow$ processamento no servidor
  $\rightarrow$ resposta com status e conteúdo $\rightarrow$ renderização e
  execução de JS $\rightarrow$ (se necessário) novas requisições via fetch.
\end{exercicios}

\section{Soluções comentadas — Capítulo 2}

\begin{exercicios}
  \item \textbf{\texttt{div} vs \texttt{section}}: \texttt{section} tem
  significado (agrupa conteúdo relacionado, idealmente com heading); \texttt{div}
  é genérico. Ex.: \texttt{<section>} para ``ingredientes''; \texttt{<div>}
  para um wrapper de layout sem significado.
  \item \textbf{Um \texttt{main} e um \texttt{h1}} por página: são os marcos
  que leitores de tela e buscadores usam para estruturar a navegação.
  \item \textbf{Formulário de login}:
\end{exercicios}

\begin{lstlisting}[style=livro, language=HTML, caption={Solução: formulário de login acessível.}, label={list:sol-login}]
<form action="/api/login" method="post">
  <div>
    <label for="email">E-mail</label>
    <input type="email" id="email" name="email"
           required autocomplete="email" />
  </div>
  <div>
    <label for="senha">Senha</label>
    <input type="password" id="senha" name="senha"
           required autocomplete="current-password" minlength="8" />
  </div>
  <button type="submit">Entrar</button>
</form>
\end{lstlisting}

\section{Soluções comentadas — Capítulo 3}

\begin{exercicios}
  \item \textbf{Especificidade vs cascata}: a especificidade decide \emph{qual
  regra} (mais específico vence); a cascata decide \emph{qual declaração}
  quando empatadas (a última escrita). Ex.: \texttt{.botao} (classe) vence
  \texttt{button} (tag); dois \texttt{.botao} — vence o último.
  \item \textbf{Box model}: sem \texttt{border-box}, \texttt{width: 300px} +
  \texttt{padding: 30px} resulta em 360px totais (conteúdo 300 + 2$\times$30).
  Com \texttt{border-box}, os 300px incluem o padding.
  \item \textbf{Centralização com Flexbox}:
\end{exercicios}

\begin{lstlisting}[style=livro, language=CSS, caption={Solução: centralização perfeita.}, label={list:sol-flex}]
.centro {
  display: flex;
  justify-content: center;
  align-items: center;
  min-height: 100vh; /* ou height fixa do container */
}
\end{lstlisting}

\section{Soluções comentadas — Capítulo 4}

\begin{exercicios}
  \item \textbf{\texttt{const} vs \texttt{let} vs \texttt{var}}:
  \texttt{const} não reatribui; \texttt{let} reatribui com escopo de bloco;
  \texttt{var} tem escopo de função e permite redeclaração (causa bugs).
  \item \textbf{\texttt{===} vs \texttt{==}}: \texttt{==} faz coerção implícita
  (\texttt{0 == false} é \texttt{true}); \texttt{===} exige valor e tipo iguais.
  \item \textbf{map/filter/reduce}:
\end{exercicios}

\begin{lstlisting}[style=livro, language=JavaScript, caption={Solução: map/filter/reduce.}, label={list:sol-map}]
const nums = [1, 2, 3, 4, 5];
const dobrados = nums.map((n) => n * 2);        // [2,4,6,8,10]
const pares = nums.filter((n) => n % 2 === 0);  // [2,4]
const soma = nums.reduce((acc, n) => acc + n, 0); // 15
\end{lstlisting}

\section{Soluções dos demais capítulos}

Os gabaritos comentados dos \emph{exercícios de fixação} de todos os
capítulos 5--25 (e os das Partes I e II acima) vivem em
\texttt{codigo/--solucoes/capNN/solucoes-capNN.md}. O arquivo
\texttt{codigo/--solucoes/README.md} explica como usar e o que cada arquivo
cobre.

\begin{caixadica}
\textbf{Por que desafios e nível sênior não têm gabarito?} São exercícios
abertos, com muitas soluções válidas. A verificação deles é o código
\emph{rodando}: cada projeto-base em \texttt{codigo/capNN-projeto/} tem
critérios de aceite e testes — se a sua solução passa nos testes e cumpre os
critérios, ela está correta. Essa é a mesma régua usada no mercado (code
review + testes), e comparar com um gabarito fechado ensinaria menos.
\end{caixadica}
